\slajdid{09_3-s004}
\begin{frame}[label=09_3-s004]
\gusttekst\setlength{\abovedisplayskip}{3pt}\setlength{\belowdisplayskip}{3pt}\setlength{\abovedisplayshortskip}{2pt}\setlength{\belowdisplayshortskip}{2pt}\begin{center}\begingroup\renewcommand{\arraystretch}{1.65}\setlength{\tabcolsep}{7pt}
\begin{tabular}{lccccc}\toprule
Režim rada & $C_{GCB}$ & $C_{GCS}$ & $C_{GCD}$ & $C_{GC}$ & $C_G$\\\midrule
zakočen & $C_{ox}WL$ & & & $C_{ox}WL$ & $C_{ox}WL+2C_0W$\\
omska oblast & 0 & $C_{ox}WL/2$ & $C_{ox}WL/2$ & $C_{ox}WL$ & $C_{ox}WL+2C_0W$\\
zasićen & 0 & $(2/3)C_{ox}WL$ & & $(2/3)C_{ox}WL$ & $(2/3)C_{ox}WL+2C_0W$\\\bottomrule
\end{tabular}\endgroup
\end{center}
Ove kapacitivnosti zavise od dimenzija tranzistora, tehnoloških parametara ali i od režima u kojem tranzistor radi. Na primer za ukupnu kapacitivnost gejta $C_G$
\par\medskip
$W$ – širina kanala\par
$L$ – dužina kanal\par
$C_{GCB}$, $C_{GCS}$, $C_{GCD}$ – kapacitivnosti gejta koje potiču od kanala, prema osnovi, sorsu, drejnu, respektivno\par
$C_{GC}$ - ukupna kapacitivnost gejta koja potiče od kanala\par
$C_0$ – jedinična kapacitivnost (normirana po širini kanala) koja potiče od preklapanja oblasti gejta i sorsa, odnosno gejta i drejna.
\end{frame}
